\documentclass{article}

\usepackage[utf8]{inputenc}
\usepackage[german]{babel}
\usepackage{listings}
\usepackage{xcolor}

\lstset{
    basicstyle=\ttfamily\small,
    breaklines=true,
    backgroundcolor=\color{gray!10}
}

\begin{document}
    \begin{center}
        \Large \textbf{Huntsman Challenge 3: Zentrales Backup- und Disaster Recovery Management} \\ [6pt]
        \normalsize
        Autor: Maximilian Jakob Maag B.Sc. \\
        Version: 1.0
    \end{center}

    \section{Intro}
    Diese Challenge behandelt die zentrale Verwaltung von Backups und Disaster Recovery Prozessen.
    Sie werden ein robustes Backup System aufbauen, das alle vom Rabbit deployten Server schützt.
    Ziel ist es, im Notfall die Infrastruktur schnell wiederherstellen zu können.

    \section{Aufgaben}

    \subsection{Aufgabe 1}
    Der Rabbit hat Backup Playbooks erstellt, die regelmäßig Datenbanken und Applikationsdaten in S3 Buckets ablegen.
    Erstellen Sie ein zentrales Inventory all dieser Backups.
    \\ \\
    Ihr Playbook sollte:
    \begin{itemize}
        \item Alle S3 Buckets aufzählen und indexieren
        \item Backup Metadaten erfassen (Datum, Größe, Integrität)
        \item Alte und verwaiste Backups identifizieren
        \item Backup Status in einer Datenbank speichern
    \end{itemize}

    \subsection{Aufgabe 2}
    Erstellen Sie ein Playbook zur Überprüfung der Backup-Integrität.
    \\ \\
    Das Playbook sollte:
    \begin{itemize}
        \item Checksums von Backup Files überprüfen
        \item Backup Archives auf Beschädigungen prüfen
        \item Stichprobenartig Backups testen (Test Restore auf temporärer VM)
        \item Bei Problemen Alerts generieren
    \end{itemize}

    \subsection{Aufgabe 3}
    Entwerfen Sie ein Backup Retention Policy Playbook für Ihre gesamte Infrastruktur.
    \\ \\
    Das Playbook sollte:
    \begin{itemize}
        \item Täglich Backups für 7 Tage speichern
        \item Wöchentliche Backups für 12 Wochen speichern
        \item Monatliche Backups für 12 Monate speichern
        \item Alte Backups automatisch löschen
        \item Replicas in mehreren Regionen halten
    \end{itemize}

    \subsection{Aufgabe 4}
    Erstellen Sie ein Disaster Recovery Testplaybook.
    \\ \\
    Dieses Playbook sollte:
    \begin{itemize}
        \item Regelmäßig (z.B. monatlich) einen Test Restore durchführen
        \item Eine Test VM mit den Backup Daten starten
        \item Die Funktionalität der wiederhergestellten Services verifizieren
        \item Ein Bericht über den DR Test Status erstellen
    \end{itemize}

    \subsection{Aufgabe 5}
    Implementieren Sie ein automated Recovery Playbook für verschiedene Fehlerszenarien:
    \begin{itemize}
        \item Service Crashes - Service automatisch neustarten
        \item Datenbank Corruption - Automatische Wiederherstellung aus letztem bekannten Good State
        \item Disk Full - Alte Log Files löschen und Space freigeben
        \item Security Breach - System isolieren und aus Backup restore
    \end{itemize}

    \subsection{Aufgabe 6}
    Erstellen Sie ein zentrales Backup Monitoring Dashboard.
    \\ \\
    Das Dashboard sollte zeigen:
    \begin{itemize}
        \item Status aller Backups (Erfolg/Fehler)
        \item Letzte Backup Zeit für jeden Host
        \item Backup Größen und Speicherkosten
        \item Recovery Point Objective (RPO) für jeden Server
        \item Time to Recover (RTO) basierend auf Testläufen
    \end{itemize}

    \subsection{Aufgabe 7}
    Dokumentieren Sie Ihren Backup- und Recovery Prozess.
    \\ \\
    Die Dokumentation sollte enthalten:
    \begin{itemize}
        \item Backup Architektur und Topologie
        \item Recovery Procedures für verschiedene Szenarien
        \item Schritt-für-Schritt Wiederherstellungsanleitung
        \item Kontakt- und Eskalationsprozeduren
    \end{itemize}

\end{document}
